\chapter{de Rham上同调与计算}
\label{ch:09-cohomology}
\ChapterPlan{周2.6.3（本地页码待核对）；梁上册第二版5.1选读5-1-2及第5章习题7；Mayer--Vietoris与球面计算交叉参考Gualtieri讲义4.3--4.4}{正文及证明初稿；附8题解答及原创图，尚未独立数学审阅}

闭形式都有局部原形式，却可能无法拼成一个全局原形式。de Rham上同调把这种障碍组成向量空间：先忽略相差一个恰当形式的区别，再用拉回、同伦和开覆盖计算这些空间。本章默认无边的有限维光滑流形，系数均为实数，使用所有光滑形式而非仅紧支形式。积分只作为检验某些类非零的工具；一般de Rham定理、奇异上同调及Poincare对偶不作为先修，也不在本章声称已证明。

\Needspace{17\baselineskip}
\section{闭形式、恰当形式与上同调}
\begin{definition}[de Rham复形与上同调]\label{def:de-rham}
第7章的$\dd^2=0$使
\[
 0\longrightarrow\Omega^0(M)\xrightarrow{\dd}\Omega^1(M)
 \xrightarrow{\dd}\cdots\xrightarrow{\dd}\Omega^n(M)\longrightarrow0
\]
成为上链复形（cochain complex）：相邻两个箭头复合为零。记闭形式空间与恰当形式空间分别为
\[
 Z^k(M)=\ker(\dd:\Omega^k\to\Omega^{k+1}),\qquad
 B^k(M)=\operatorname{im}(\dd:\Omega^{k-1}\to\Omega^k).
\]
由于$B^k\subset Z^k$，定义de Rham上同调（de Rham cohomology）
\[
 H^k_{\mathrm{dR}}(M)=Z^k(M)/B^k(M).
\]
闭形式$\omega$的类记为$[\omega]$；两个闭形式代表同一类，当且仅当相差$\dd\alpha$。约定负次数形式为零，所以$B^0=0$；$k<0$或$k>n$时$H^k_{\mathrm{dR}}(M)=0$。下文简记$H^k(M)$，始终指普通实de Rham上同调。
\end{definition}
商空间中的零类意味着存在一个\emph{全局}原形式。所有形式局部恰当并不使这个商为零；局部原形式之间的差异必须能够相容消除，才会给出全局原形式。

\begin{proposition}[零次上同调与分支]\label{prop:cohomology-zero}
$H^0(M)$是局部常值函数的空间。若$\pi_0(M)$为连通分支集合，则
\[
 H^0(M)\simeq\prod_{C\in\pi_0(M)}\R.
\]
特别地，非空连通流形的$H^0$为$\R$。
\end{proposition}
\begin{proof}
$\dd f=0$说明小坐标球内所有偏导为零。沿球内线段积分得到函数恒定，所以$f$局部常值。反过来局部常值函数的微分为零。流形局部道路连通，故连通分支为开集；局部常值函数在每个连通分支上恒定，而任意逐分支指定的实数给出光滑局部常值函数。$B^0=0$，无需再取非平凡的商。若分支无限多，每个分支上都可非零，所以是直积而非有限支撑直和。
\end{proof}
\begin{example}[维数不能从一张卡推断]
可数离散零维流形$M=\{p_1,p_2,\ldots\}$的$H^0(M)$是所有实数列，已经无限维。对非空星形域$U\subset\R^n$，第7章Poincare引理则给$H^0(U)=\R$、$H^k(U)=0$（$k>0$）。一般流形即使每张小卡都如此，全局上同调也不必消失。
\end{example}

\begin{proposition}[外积下降到上同调]\label{prop:cohomology-product}
若$\alpha\in Z^p(M),\beta\in Z^q(M)$，则
$[\alpha]\,[\beta]=[\alpha\wedge\beta]$良定义，使$H^\bullet(M)$成为分次交换实代数，满足$ab=(-1)^{pq}ba$。单位为常函数$1$的类（$M$非空）。
\end{proposition}
\begin{proof}
乘法法则给$\dd(\alpha\wedge\beta)=0$。若将$\alpha$改成$\alpha+\dd\xi$，乘积改变为$\dd\xi\wedge\beta=\dd(\xi\wedge\beta)$。若将$\beta$改成$\beta+\dd\zeta$，改变为$\alpha\wedge\dd\zeta=(-1)^p\dd(\alpha\wedge\zeta)$。依次改变两项，类仍不变。结合性、分配律、单位和分次交换律均由形式外积继承。因此奇数次数类$a$满足$a^2=-a^2$，实系数下推出$a^2=0$；偶数次数时不能据此推断平方为零。
\end{proof}

\begin{proposition}[拉回的反变性]\label{prop:cohomology-pullback}
任意光滑$F:M\to N$诱导保次数的代数同态$F^*:H^\bullet(N)\to H^\bullet(M)$，定义为$F^*[\omega]=[F^*\omega]$，且$(G\circ F)^*=F^*\circ G^*$、$\operatorname{id}^*=\operatorname{id}$。微分同胚诱导同构。
\end{proposition}
\begin{proof}
自然性$\dd F^*=F^*\dd$保证闭形式仍闭，且$F^*(\dd\alpha)=\dd(F^*\alpha)$，所以商空间上的映射与代表元无关。拉回保持外积；复合与恒等关系在形式层成立，故传到类上。若$F$有光滑逆，则其逆拉回恰好给$F^*$的逆。无需$F$单射、满射或满秩才能定义拉回。
\end{proof}

\Needspace{12\baselineskip}
\begin{proposition}[周期检验与顶次积分]\label{prop:cohomology-period}
设$S$为紧、无边、定向的$k$维流形，$k\geq1$，$f:S\to M$光滑。则$[\omega]\mapsto\int_S f^*\omega$是$H^k(M)$上的良定义线性泛函。特别地，紧、无边、定向$n$维$M$上的顶次积分下降为$H^n(M)\to\R$。
\end{proposition}
\begin{proof}
若$\omega$改为$\omega+\dd\alpha$，积分之差为$\int_S\dd(f^*\alpha)$，由第8章Stokes及$\partial S=\varnothing$等于零。$S$紧使原形式紧支，无需额外支集条件。线性来自积分线性。
\end{proof}
对一般紧连通定向流形，这个泛函非零：正体积形式给正积分。但此时尚未证明它必为同构，也未证明所有类均能由周期区分。后面会对球面用维数计算单独证明同构，不能以一般de Rham定理替代本章缺失的步骤。
\begin{example}[穿孔平面上的非零类]
第7章的$\eta=(-y\dd x+x\dd y)/(x^2+y^2)$闭，逆时针圆周周期为$2\pi$，所以$[\eta]\ne0$。周期不为零能够排除恰当；一次周期恰好为零通常不足以判定任意流形上的类为零。对圆周，单个基本周期就足够，须由后面的完整计算证明。
\end{example}


\Needspace{15\baselineskip}
\section{同伦不变性}
\begin{definition}[光滑同伦与光滑同伦等价]
光滑映射$F_0,F_1:M\to N$之间的光滑同伦（smooth homotopy）是$H:[0,1]\times M\to N$，在$[0,1]\times M$的某开邻域可光滑延拓，且$H(0,p)=F_0(p),H(1,p)=F_1(p)$。若$f:M\to N,g:N\to M$满足$g\circ f$光滑同伦于$\operatorname{id}_M$、$f\circ g$光滑同伦于$\operatorname{id}_N$，称$f,g$为光滑同伦等价。
\end{definition}
这里明确使用光滑同伦；从连续同伦推出光滑同伦还需逼近定理，本章不提前调用。所用的收缩和形变均会给出光滑公式及定义域。

\begin{theorem}[光滑同伦不变性]\label{thm:cohomology-homotopy}
若$F_0,F_1:M\to N$光滑同伦，则它们在所有次数上诱导相同的拉回。光滑同伦等价诱导de Rham上同调代数同构。
\end{theorem}
\begin{proof}
由第7章定理\ref{thm:homotopy-operator}，$F_1^*\omega-F_0^*\omega=\dd K_H\omega+K_H\dd\omega$。对闭$k$形式，$k\geq1$时差为恰当形式$\dd K_H\omega$，所以两类相同。$k=0$时$K_H\omega=0$且$\dd\omega=0$，同一公式直接给$F_1^*\omega=F_0^*\omega$，没有负次数原形式。
若$f,g$为光滑同伦等价，反变复合律给$f^*g^*=(g\circ f)^*=\operatorname{id}$和$g^*f^*=\operatorname{id}$，所以二者互逆，且保持外积。
\end{proof}

\begin{corollary}[光滑形变收缩]\label{cor:deformation-retract}
设$i:A\hookrightarrow M$为子流形包含，$r:M\to A$光滑，$r\circ i=\operatorname{id}_A$。若$i\circ r$光滑同伦于$\operatorname{id}_M$，则$i^*:H^\bullet(M)\to H^\bullet(A)$与$r^*$互逆。若$M$光滑收缩到一点，则$H^0(M)=\R$且正次数上同调为零。
\end{corollary}
\begin{proof}
首项直接用上一条定理。收缩到一点时取$A=\{p\}$；一点上的零形式为实数，正次数形式全为零，所以得到所列空间。此处未要求$M$具有全局星形坐标；收缩的光滑性才是使用的假设。
\end{proof}

\Needspace{20\baselineskip}
\begin{example}[乘实线与径向形变]\label{ex:radial-retract}
投影$\pi:M\times\R\to M$、零截面$i(p)=(p,0)$满足$\pi i=\operatorname{id}$，$H_t(p,s)=(p,(1-t)s)$将恒等映射光滑变到$i\pi$。因此$H^\bullet(M\times\R)\simeq H^\bullet(M)$。同理可处理$\R^m$因子。

在$\R^{n+1}\setminus\{0\}$，令$r(x)=x/|x|$、$i:S^n\hookrightarrow\R^{n+1}\setminus\{0\}$。公式
\[
 H_t(x)=\bigl((1-t)+t/|x|\bigr)x
\]
给$\operatorname{id}$到$ir$的光滑同伦：$0\leq t\leq1$时括号严格正，轨迹不经过原点。由括号保持正的开集，得到所需开邻域延拓；不需要对所有$x$统一延长时间。故穿孔空间与单位球面的普通de Rham上同调同构。$n=0$时单位球面为两点，这一例子仍适用，但不能把它当成连通球面。
\end{example}
\begin{figure}[htbp]
\centering\DGDiagram{radial-retract}

\caption{穿孔平面到单位圆的径向形变：内外两点都沿不经过原点的正半径移动到圆上。箭头表示$t$从$0$到$1$的方向；空圈为删去的原点。原创示意。}
\label{fig:radial-retract}
\end{figure}

\begin{example}[紧支条件不能自动传递]
常值映射$\R\to\{p\}$把一点上非零的紧支函数拉回为实线上的非紧支常函数。一般光滑拉回不一定保持紧支，同伦算子也不一定保持紧支；本节的普通上同调同伦不变性不能直接套用于紧支上同调。习题9.8将计算$H_c^1(\R)$，并与$H^1(\R)=0$比较。
\end{example}


\Needspace{15\baselineskip}
\section{Mayer--Vietoris序列}
局部原形式在交叠上的不相容，是把局部计算拼成全局计算时要保存的数据。Mayer--Vietoris序列把全域、两个开块及其交的上同调连成一个正合序列；连接映射恰好记录这种拼接障碍。
\begin{definition}[正合序列与本节符号]
线性映射序列$A\xrightarrow{f}B\xrightarrow{g}C$在$B$处正合（exact），指$\operatorname{im}f=\ker g$。设$M=U\cup V$为两个开集的并，$W=U\cap V$，规定
\[
 r(\omega)=(\omega|_U,\omega|_V),\qquad
 q(\alpha,\beta)=\alpha|_W-\beta|_W.
\]
后面的连接映射$\delta$始终对应这个差的顺序。改变差的顺序会改变连接映射的号，不能在公式和计算间混用。
\end{definition}

\begin{lemma}[形式复形的短正合序列]\label{lem:mv-short}
对每个次数$k$，有短正合序列
\[
 0\longrightarrow\Omega^k(M)\xrightarrow{r}
 \Omega^k(U)\oplus\Omega^k(V)\xrightarrow{q}\Omega^k(W)
 \longrightarrow0,
\]
且$r,q$与$\dd$交换。
\end{lemma}
\begin{proof}
形式在开覆盖上由限制唯一确定，故$r$单射；$q(\alpha,\beta)=0$恰好说两者在交叠相同，由形式的局部拼接得到唯一全局形式，所以$\ker q=\operatorname{im}r$。

证明$q$满射需要第8章的单位分解。取$\rho_U+\rho_V=1$，$\operatorname{supp}\rho_U\subset U$、$\operatorname{supp}\rho_V\subset V$。若先得到多项从属分解，将指定给同一开集的各项相加；局部有限性保证和光滑，所涉闭支集的局部有限并仍闭，因此和的支集不越出对应开集。这里的两函数未必紧支，也无此要求。

对$\gamma\in\Omega^k(W)$，令
\[
 a=(\rho_V\gamma)^{\mathrm{ext}}\in\Omega^k(U),\qquad
 b=(-\rho_U\gamma)^{\mathrm{ext}}\in\Omega^k(V),
\]
其中上标表示在各自定义域内、交叠之外延零。例如$p\in U\setminus V$时，$p\notin\operatorname{supp}\rho_V$，邻域中$\rho_V$恒为零，故$a$延零光滑，即使$\gamma$在交叠边缘无界也无妨。$b$同理。在$W$上$a-b=(\rho_V+\rho_U)\gamma=\gamma$，故$q$满射。限制与外微分交换，所以$r,q$均为上链映射。
\end{proof}
注意：满射给出的$a,b$通常不闭。若它们总能闭，就会抹掉所需的连接映射以及非平凡的全局上同调。

\begin{proposition}[连接映射的构造]\label{prop:mv-connecting}
给定闭$k$形式$\gamma$于$W$，取任意提升$q(a,b)=\gamma$。则$\dd a,\dd b$可拼成全局闭$(k+1)$形式$\eta$，并可定义
\[
 \delta:H^k(W)\to H^{k+1}(M),\qquad\delta[\gamma]=[\eta].
\]
它与提升、代表元及构造提升所用单位分解无关。用上一引理的提升时，
\[
 \eta|_U=\dd(\rho_V\gamma)^{\mathrm{ext}},\qquad
 \eta|_V=-\dd(\rho_U\gamma)^{\mathrm{ext}}.
\]
在$W$上二者共同等于$\dd\rho_V\wedge\gamma=-\dd\rho_U\wedge\gamma$。
\end{proposition}
\begin{proof}
$\dd\gamma=0$给$\dd a|_W-\dd b|_W=0$，所以两者相容并唯一拼成$\eta$。其微分在$U,V$均为零，故全局为零。

若换提升为$(a',b')$，差在$\ker q$中，故由短正合性为$r(\xi)$，$\xi\in\Omega^k(M)$。于是新拼接形式与$\eta$的差为$\dd\xi$，上同调类不变。若将$\gamma$改为$\gamma+\dd\lambda$（$k\geq1$），选$q(c,d)=\lambda$，则$(a+\dd c,b+\dd d)$提升新代表元，外微分仍为$(\dd a,\dd b)$，故类不变。$k=0$时没有负次数的$\lambda$，只需前一种检查。任一单位分解只是给出一种提升，已经涵盖在提升无关性中。提升可按线性组合相加，故$\delta$线性。
\end{proof}

\Needspace{17\baselineskip}
\begin{theorem}[Mayer--Vietoris长正合序列]\label{thm:mayer-vietoris}
对开覆盖$M=U\cup V$，上述$r,q,\delta$诱导长正合序列
\begin{align*}
 0&\longrightarrow H^0(M)\xrightarrow{r_*}H^0(U)\oplus H^0(V)
 \xrightarrow{q_*}H^0(W)\xrightarrow{\delta}H^1(M)\longrightarrow\cdots,\\
 \cdots&\longrightarrow H^k(M)\xrightarrow{r_*}H^k(U)\oplus H^k(V)
 \xrightarrow{q_*}H^k(W)\xrightarrow{\delta}H^{k+1}(M)\longrightarrow\cdots.
\end{align*}
\end{theorem}
\begin{proof}
分别检查三类位置；所有形式均为全域或相应开块上的光滑形式，不加紧支限制。

\textbf{在$H^k(U)\oplus H^k(V)$处。}来自同一全局闭形式的限制相同，故$\operatorname{im}r_*\subset\ker q_*$。反向，设$\alpha,\beta$闭且$q_*([\alpha],[\beta])=0$。当$k\geq1$，有$\alpha|_W-\beta|_W=\dd\gamma$。取$q(a,b)=\gamma$，则$\alpha-\dd a,\beta-\dd b$闭且在$W$上相同，可拼成全局闭形式$\omega$；它限制后的类分别是$[\alpha],[\beta]$。当$k=0$，零类只能是零函数，故$\alpha|_W=\beta|_W$，直接拼接即可。

\textbf{在$H^k(W)$处。}若$[\gamma]=q_*([\alpha],[\beta])$且$\alpha,\beta$闭，以$(\alpha,\beta)$为提升即可得到$\delta[\gamma]=0$。反向，取$q(a,b)=\gamma$并将$\dd a,\dd b$拼成$\eta$。若$\delta[\gamma]=0$，则$\eta=\dd\xi$，$\xi\in\Omega^k(M)$。令$\alpha=a-\xi|_U,\beta=b-\xi|_V$，它们闭且$q(\alpha,\beta)=\gamma$。所以$\ker\delta\subset\operatorname{im}q_*$。

\textbf{在$H^{k+1}(M)$处。}$\delta[\gamma]$的代表元在$U,V$分别为$\dd a,\dd b$，故限制类均为零，得到$\operatorname{im}\delta\subset\ker r_*$。反向，若全局闭形式$\eta$满足$\eta|_U=\dd a,\eta|_V=\dd b$，则$\gamma=a|_W-b|_W$闭；以$(a,b)$为提升得到$\delta[\gamma]=[\eta]$。这证明所有正次数的该位置正合。最左端$H^0(M)$的单射来自局部常值函数由覆盖上的限制唯一确定，补足序列起点。
\end{proof}
\begin{figure}[htbp]
\centering\DGDiagram{mv-connecting}

\caption{连接映射的构造：先将交叠闭形式$\gamma$提升为不必闭的$(a,b)$，再微分；$\dd a,\dd b$相容，故来自全局闭形式$\eta$。箭头$r$为限制，不把它画成任意局部数据的全局延拓。原创示意。}
\label{fig:mv-connecting}
\end{figure}

\begin{example}[为什么圆周的交叠必须记住两个分支]
取$U=S^1\setminus\{(-1,0)\}$、$V=S^1\setminus\{(1,0)\}$。两块都是开弧，但$W$是上下两个不相交开弧，故$H^0(W)=\R^2$。若按“每块可收缩”把交叠也误记为一个点，就会丢掉一维余核。具体矩阵及连接映射的积分符号将在下一节和习题9.5中核对。
\end{example}
本节没有假定上同调有限维，也不要求开覆盖成员紧致。正合性只给核与像的关系，不能把任意长正合序列直接拆成各项的无条件直和分解。


\Needspace{15\baselineskip}
\section{圆周与球面的上同调计算}
\begin{theorem}[圆周的周期与上同调]\label{thm:circle-cohomology}
对逆时针定向的单位圆，记$\lambda=(x\dd y-y\dd x)|_{S^1}$。则
\[
 H^k(S^1)=\begin{cases}\R,&k=0,1,\\0,&\text{其余次数},\end{cases}
 \qquad [\alpha]\longmapsto\frac1{2\pi}\int_{S^1}\alpha
\]
是$H^1(S^1)$到$\R$的同构，将$[\lambda]$送为$1$。其上同调代数由单位和一次元$e=[\lambda]$生成，$e^2=0$。
\end{theorem}
\begin{proof}
圆周连通，故$H^0=\R$，高于一次的形式为零。每个一形式因维数为一均闭。令$\pi:\R\to S^1$，$\pi(t)=(\cos t,\sin t)$。这是局部微分同胚，$\pi^*\alpha=f(t)\dd t$且$f$光滑、$2\pi$周期。又$\pi^*\lambda=\dd t$。
取$a=(2\pi)^{-1}\int_0^{2\pi}f(t)\dd t$，令
\[
 h(t)=\int_0^t(f(s)-a)\dd s.
\]
由于$f-a$在任一长度$2\pi$区间上的积分为零，$h(t+2\pi)=h(t)$。故$h$下降为$S^1$上的光滑函数$g$；光滑性可在$\pi$的局部逆上检验。$\pi^*(\alpha-a\lambda)=\dd h=\pi^*\dd g$，而满射局部微分同胚的拉回在形式层单射，所以$\alpha=a\lambda+\dd g$。这证明每个类由一个$a$确定。反向，$\int\dd g=0$而$\int\lambda=2\pi$，所以不同$a$给不同类。平方因二形式为零而为零，得到代数描述。
\end{proof}

\begin{example}[用Mayer--Vietoris复核圆周]
沿上一节的两开弧覆盖，以$W$的上、下分支为顺序，长序列缩为
\[
 0\longrightarrow\R\xrightarrow{c\mapsto(c,c)}\R^2
 \xrightarrow{(a,b)\mapsto(a-b,a-b)}\R^2
 \xrightarrow{\delta}H^1(S^1)\longrightarrow0.
\]
中间映射矩阵为$\left(\begin{smallmatrix}1&-1\\1&-1\end{smallmatrix}\right)$，秩为一，像是对角线。因此$H^1(S^1)\simeq\R^2/\{(c,c)\}$，维数一。与周期定理比较时还须固定符号：按本章$q=\alpha-\beta$，习题9.5证明交叠常数$(c_+,c_-)$满足$\int_{S^1}\delta(c_+,c_-)=c_+ -c_-$。仅维数相同并不能自动确定某个连接类对应$+\lambda$还是$-\lambda$。
\end{example}

\begin{lemma}[球面两极覆盖的交叠]\label{lem:sphere-overlap}
对$n\geq1$，令$N=(0,\ldots,0,1),S=-N$，$U=S^n\setminus\{S\}$、$V=S^n\setminus\{N\}$。两块各微分同胚于$\R^n$，而$W=U\cap V$光滑形变收缩到赤道$S^{n-1}$。
\end{lemma}
\begin{proof}
立体投影及其光滑逆已在第3章构造，所以$U,V$的正次数上同调为零。对$W$内的点写$x=(x',z)$，$-1<z<1$，$x'\ne0$，则
\[
 x=\bigl(\sqrt{1-z^2}\,u,z\bigr),\qquad u=x'/|x'|\in S^{n-1}.
\]
该式及其逆给$W\simeq S^{n-1}\times(-1,1)$的微分同胚。保持$u$不变，将$z$变为$(1-t)z$，即得赤道收缩。在满足$|(1-t)z|<1$的开集中公式仍光滑，故有本章要求的开邻域延拓；$t\in[0,1]$时这一不等式始终成立。$n=1$时$u\in S^0=\{-1,1\}$，交叠正好是两个分支，不能合并。
\end{proof}
\begin{figure}[htbp]
\centering\DGDiagram{sphere-mv-cover}

\caption{球面的两极开覆盖及交叠收缩。图示为二维球面：$U$删去南极，$V$删去北极，交叠含赤道且沿纬度收缩到赤道；不是把闭半球当作开覆盖。橙色箭头为收缩方向。原创示意。}
\label{fig:sphere-mv-cover}
\end{figure}

\begin{theorem}[所有正维球面的计算]\label{thm:sphere-cohomology}
对$n\geq1$，
\[
 H^k(S^n)=\begin{cases}\R,&k=0\text{或}k=n,\\0,&\text{其余次数}.\end{cases}
\]
若$S^n$取定向，积分$H^n(S^n)\to\R$是同构。取积分为$1$的顶次类$u$，其上同调代数由$1,u$生成且$u^2=0$，$\deg u=n$。$S^0$须另计：$H^0(S^0)=\R^2$，正次数全为零。
\end{theorem}
\begin{proof}
$S^1$已证明。设$n\geq2$，使用上引理的覆盖。$U,V$连通且正次数上同调为零，$W$与连通的$S^{n-1}$同伦等价。长序列的低次部分是
\[
 0\longrightarrow H^0(S^n)\longrightarrow\R^2
 \xrightarrow{(a,b)\mapsto a-b}\R\longrightarrow H^1(S^n)
 \longrightarrow0.
\]
减法映射满射、核为对角线，故$H^0(S^n)=\R$、$H^1(S^n)=0$。对$k\geq2$，相应一段为
\[
 0\longrightarrow H^{k-1}(W)\xrightarrow{\delta}H^k(S^n)
 \longrightarrow0,
\]
因为两侧来自$U,V$的$k-1$及$k$次群均为零。因此$H^k(S^n)\simeq H^{k-1}(S^{n-1})$。以$n=1$为起点归纳：只有$k=n$留下$\R$，其余$k\geq2$为零；高于维数也由形式为零直接验证。这一归纳没有在$k=1$处误用正次数消失，因为$H^0(U),H^0(V)$并不为零。

取$\R^{n+1}$标准体积$\nu$、位置向量场$R=\sum_{i=1}^{n+1}x^i\partial_i$，$j:S^n\hookrightarrow\R^{n+1}$。$\sigma=j^*(\iota_R\nu)$是外向定向的正体积形式：球面上$R$就是外单位法向，在外向首位正基上求值为正。因球面紧、无边，积分在类上良定义，且$A=\int_{S^n}\sigma>0$。由已证明的$\dim H^n(S^n)=1$，非零线性泛函$\int$必为同构。取$u=[\sigma/A]$，得到积分为$1$的生成元。若反转定向，改取相应负号即可。$u^2$因$2n>n$为零，与$n$奇偶无关。$S^0$为两点，直接用零次分支命题计算；其乘法是$\R^2$上的逐分量乘法。
\end{proof}

\begin{example}[穿孔空间与环域的结果]
对$m\geq2$，径向收缩和球面计算给$\R^m\setminus\{0\}$的上同调仅在$k=0,m-1$为$\R$。$m=2$时角形式生成$H^1$；$m\geq3$时$H^1=0$，所以每个闭一形式均有全局原函数，虽该域不是星形域。$m=1$时域为两条开半线，$H^0=\R^2$，正次数为零。平面圆环同样径向收缩到一条圆周，故闭一形式类由绕孔的一个周期确定。
\end{example}
圆周与球面只展示了几种完整计算；这里未由此声称所有流形上同调有限维，也未建立一般Kunneth公式、Poincare对偶或与奇异上同调的比较定理。


\Needspace{15\baselineskip}
\section{上同调习题及详细解答}
习题9.1以梁书上册第二版第5章原题7（印刷页131/PDF147）为出发点，将原题的拓扑同胚条件加强为全局坐标微分同胚，并构造原函数；不据此声称已证明仅有拓扑同胚时的完整原题。习题9.5以Gualtieri公开讲义4.4例4.18后的留证任务为出发点，按本书差映射重写并核对符号；无梁书原题号。其余6题为原创补充，全部解答由本工程编写。

来源与题号映射见\texttt{references/EXERCISES.md}。

\Needspace{12\baselineskip}
\subsection*{习题9.1：闭一形式与全局坐标}
采用梁书第5章原题7的加强假设版本。设流形开集$O$有全局坐标微分同胚$x:O\to\R^n$，$\alpha=\alpha_i\dd x^i$。证明$\partial_j\alpha_i=\partial_i\alpha_j$当且仅当存在$f\in C^\infty(O)$使$\alpha=\dd f$，并给出$f$的积分公式。若坐标像只是任意欧氏开集，结论是否仍成立？
\begin{proof}[解答]
若$\alpha=\dd f$，混合偏导交换给出条件。反向，所列条件等价于$\dd\alpha=0$。令$\widetilde\alpha=(x^{-1})^*\alpha$，在星形域$\R^n$上取
\[
 \widetilde f(z)=\int_0^1\widetilde\alpha_{tz}(z)\dd t
 =\int_0^1\alpha_i(x^{-1}(tz))z^i\dd t.
\]
第7章径向公式给$\dd\widetilde f=\widetilde\alpha$，故$f=\widetilde f\circ x$满足要求。这里$x^{-1}(0)$存在，整条线段均在坐标像中。若坐标像为穿孔平面，则角形式闭却不恰当，结论失败。仅拓扑同胚而非光滑坐标微分同胚也不允许在其下直接变换光滑形式；本题实际使用的是光滑坐标以及坐标像全为$\R^n$。
\end{proof}

\Needspace{12\baselineskip}
\subsection*{习题9.2：代表元与外积}
设$\alpha,\beta$分别为$p,q$次闭形式，$p,q\geq1$，$\xi,\zeta$次数分别为$p-1,q-1$。直接写出$(\alpha+\dd\xi)\wedge(\beta+\dd\zeta)-\alpha\wedge\beta$的一个原形式，并说明这如何证明上同调乘积良定义。
\begin{proof}[解答]
展开差为$\dd\xi\wedge\beta+\alpha\wedge\dd\zeta+\dd\xi\wedge\dd\zeta$。令
\[
 \Theta=\xi\wedge\beta+(-1)^p\alpha\wedge\zeta+\xi\wedge\dd\zeta.
\]
由$\dd\alpha=\dd\beta=\dd^2\zeta=0$，乘法法则逐项给$\dd\Theta$等于上述三项之和。特别是第二项中的$(-1)^p$抵消跨过$p$次形式产生的号；第三项不再需要额外系数。因此更换两个闭代表元只改变一个恰当形式，不改变外积类。这个结论不要求$\alpha,\beta$自身恰当。
\end{proof}

\Needspace{12\baselineskip}
\subsection*{习题9.3：圆环上的完整判别}
设$A=\{(x,y):a<\sqrt{x^2+y^2}<b\}$，$0<a<b$，$\alpha$为闭一形式。取任意$r_0\in(a,b)$的逆时针圆周$C_{r_0}$。证明$\alpha$恰当当且仅当$\int_{C_{r_0}}\alpha=0$，并说明为何一个周期在这里足够。
\begin{proof}[解答]
径向映射$r(x)=r_0x/|x|$给$A$到$C_{r_0}$的收缩，线性插值的半径$(1-t)|x|+tr_0$仍位于$(a,b)$，所以是光滑形变收缩。于是包含$i:C_{r_0}\hookrightarrow A$诱导同构$i^*:H^1(A)\to H^1(C_{r_0})$。缩放圆周到单位圆，再用定理\ref{thm:circle-cohomology}，知$[i^*\alpha]=0$恰好等价于所列周期为零。由$i^*$单射，这又等价于$[\alpha]=0$，即有全局原函数。必要性也可直接由全微分的端点公式得到。一个周期足够是因为已经证明这里$H^1$一维且周期给同构，并非任意流形都可照此判断。
\end{proof}

\Needspace{12\baselineskip}
\subsection*{习题9.4：绕圈映射的拉回}
对整数$m$，令$F_m:S^1\to S^1$满足$F_m(e^{it})=e^{imt}$。计算$F_m^*$在$H^0,H^1$上的作用，分别解释$m=0$及$m<0$。若$m\ne1$，能否有到恒等映射的光滑同伦？
\begin{proof}[解答]
整数性使$t$换成$t+2\pi$不改变映射，所以$F_m$良定义；正数幂、零次常值及负数幂在单位圆上均光滑。对圆周形式$\lambda$，用参数拉回得到$\pi^*F_m^*\lambda=m\dd t=\pi^*(m\lambda)$，故$F_m^*\lambda=m\lambda$。$H^1$上为实数乘$m$，$H^0$上的常函数拉回仍为自身，即恒等。$m=0$是常值映射，杀掉一次类却不杀掉零次常数；$m<0$使一次周期按负数$m$变化。若$m\ne1$且有到恒等映射的光滑同伦，同伦不变性会迫使其$H^1$拉回既为乘$m$又为恒等，矛盾。这里不必先建立一般映射度理论。
\end{proof}

\Needspace{12\baselineskip}
\subsection*{习题9.5：连接映射的周期符号}
沿正文圆周两开弧覆盖，取$U$删去左端点、$V$删去右端点，$W_+,W_-$分别为上、下开弧。设$\gamma$在两分支上分别为$c_+,c_-$。从$q=\alpha-\beta$构造$\delta[\gamma]$的代表元，求其逆时针周期，并给出$\delta(1,0)$与$[\lambda]$的关系。
\begin{proof}[解答]
选光滑$\chi:[-1,1]\to[0,1]$，在$x=-1$附近为$1$、在$x=1$附近为$0$，令$\rho_V=\chi(x)$、$\rho_U=1-\chi(x)$，限制在圆周上。它们的支集分别避开右、左端点，满足从属性。按命题\ref{prop:mv-connecting}，在$W$上连接形式是$\eta=\gamma\dd\rho_V$，在两个被删端点附近因$\dd\rho_V=0$延为零，得到全局光滑一形式。

逆时针上弧从右端到左端，$\rho_V$由$0$变$1$；下弧从左端到右端，由$1$变$0$。基本定理给
\[
 \int_{S^1}\eta=c_+(1-0)+c_-(0-1)=c_+-c_-.
\]
因此$\delta(1,0)$周期为$1$，由圆周周期同构，它等于$[\lambda/(2\pi)]$，不是$[\lambda]$。对角常数$(c,c)$的连接类为零，与$q_*$的像吻合。若改变差映射顺序，则连接代表元及周期都应同时变号。
\end{proof}

\Needspace{12\baselineskip}
\subsection*{习题9.6：穿孔三维空间的两个次数}
求$H^1(\R^3\setminus\{0\})$和$H^2(\R^3\setminus\{0\})$。构造一个闭而不恰当的二形式，并解释为什么同一空间上的每个闭一形式却都恰当。
\begin{proof}[解答]
径向收缩$r(x)=x/|x|$给与$S^2$的同伦等价，所以$H^1=0,H^2=\R$。在穿孔空间令
\[
 \omega=\frac{x\dd y\wedge\dd z-y\dd x\wedge\dd z+z\dd x\wedge\dd y}
 {(x^2+y^2+z^2)^{3/2}}.
\]
其外微分系数是向量函数$x/|x|^3$的欧氏散度，直接计算为$3|x|^{-3}-3|x|^{-5}(x^2+y^2+z^2)=0$，所以闭。限制到单位球面就是第8章面积形式，积分$4\pi$；若$\omega=\dd\beta$，在球面拉回后Stokes应给零，矛盾。另一方面$H^1=0$正是说所有闭一形式的类为零，故都有全局原函数。域缺一个点并不意味着每个正次数都有非零上同调；这里的障碍出现在二次。
\end{proof}

\Needspace{12\baselineskip}
\subsection*{习题9.7：球面顶次类的归一化}
在外向定向的$S^2$上，设$\sigma$为面积形式，$\omega$为任意光滑二形式。证明
\[
 [\omega]=\frac1{4\pi}\left(\int_{S^2}\omega\right)[\sigma],
\]
并由此判断零积分二形式是否恰当。该判断是否已对所有紧定向流形证明？
\begin{proof}[解答]
所有二形式在二维球面上自动闭。球面定理给积分$H^2(S^2)\to\R$为同构，而$\int\sigma=4\pi$。因此$\omega-(4\pi)^{-1}(\int\omega)\sigma$的积分为零，其类必为零，得到所列等式。特别地，$\int\omega=0$当且仅当$\omega$恰当。这是存在一个全局一形式原形式的结论，不声称上式本身已经写出该原形式。对所有紧连通定向流形的一般顶次积分同构，本章尚未证明；不能将球面上的维数计算扩展为一般证明。
\end{proof}

\Needspace{12\baselineskip}
\subsection*{习题9.8：紧支与普通上同调的差别}
定义$H_c^k(M)=\ker(\dd:\Omega_c^k\to\Omega_c^{k+1})/\dd\Omega_c^{k-1}$。证明实线上的积分诱导$H_c^1(\R)\simeq\R$，而$H^1(\R)=0$。为何这不违反光滑同伦不变性？
\begin{proof}[解答]
在实线上紧支一形式均形如$f(x)\dd x$，自动闭。紧支函数$g$满足$\int_\R\dd g=g(+\infty)-g(-\infty)=0$，所以积分在$H_c^1$上良定义。若$\int_\R f=0$，令$F(x)=\int_{-\infty}^x f(s)\dd s$。它在$f$支集左侧为零，在右侧等于总积分零，所以$F$紧支且光滑，$\dd F=f\dd x$。这证明积分的核恰是紧支恰当形式。取一个非负非零光滑紧支函数并除以其正积分，可得积分为$1$的代表元，故该泛函满射，因而为同构。

普通上同调允许原函数不紧支；任意光滑$f$均有$G(x)=\int_0^x f(s)\dd s$，故$H^1(\R)=0$。实线确实光滑收缩到一点，但本章同伦不变性针对所有光滑形式的复形。一般拉回与同伦算子不保持紧支，不能把同一个论证移到$H_c$。本题只计算此例，未建立紧支上同调在适当映射下的完整理论。
\end{proof}

